% !TeX root = main.tex
\chapter{Conclusion and Future Work}
\label{ch:conclusions}

\section{Summary of Findings}
\label{sec:conclusions}

This thesis shows that, with a few annotated pages, the quality of \gls{tls} in historical manuscripts is determined by the initialization and the line reconstruction, not by the encoder architecture. The proposed models reach the state of the art, with 92.20\% on U-DIADS-TL, compared with the best published 83.00\%~\cite{Zottin_2025_FEST}, and 99.43\% on DIVA-HisDB, compared with 97.43\%~\cite{signals3030032}. Neither the family nor the size of the encoder changes the results as much as the initialization and the line reconstruction. With pretraining on annotated manuscript lines, three training pages per manuscript are sufficient for the state of the art, while the result with one page depends on the selected page. In the cross-collection experiments, the center-line model reaches higher line detection metrics on unseen collections than Mask R-CNN, with 93.07\% compared with 49.09\% \gls{fm} without fine-tuning.

\section{Answering the Research Questions}
\label{sec:disc-answers}\label{sec:disc-interpretation}\label{sec:disc-initialization}\label{sec:disc-representations}\label{sec:disc-previous}

The following paragraphs answer the three research questions based on the results of Chapter~\ref{ch:results}.

\textbf{RQ1: How do the encoder architecture, its initialization, and the line reconstruction affect few-shot text line segmentation in historical manuscripts and yield \gls{sota} results in terms of line detection metrics?}

The transformer-based encoders do not perform better than the \gls{cnn} encoders, as PVTv2-B2 and ViT-Small are better than ConvNeXt-Tiny in only 10 of 20 and 10 of 26 paired comparisons in Table~\ref{tab:rq1-paired}. The initialization has a larger and more consistent effect, as CATMuS is better than ImageNet in 11 of 12 paired comparisons by up to 20.76 percentage points, which agrees with the benefit of in-domain pretraining reported by De Nardin et al.~\cite{denardin2025indomain}. The largest effect comes from the line reconstruction, which raises the U-DIADS-TL result of the same Mask R-CNN detections from 16.61\% to 90.82\% and leads to 92.20\% on U-DIADS-TL and 99.43\% on DIVA-HisDB, compared with the best published 83.00\%~\cite{Zottin_2025_FEST} and 97.43\%~\cite{signals3030032}.

\textbf{RQ2: How many labeled pages are needed to achieve performance comparable to \gls{sota} methods in terms of line detection metrics?}

Three training pages per manuscript are sufficient if the model is pretrained on annotated lines, although additional annotated pages are still used for validation. With CATMuS initialization, Mask R-CNN reaches 98.67 \(\pm\) 0.29\% on DIVA-HisDB with three pages over three seeds, compared with 97.43\% for Mask R-CNN+ trained on all pages~\cite{signals3030032}. With one page, the result depends on the selected page and is 95.39 \(\pm\) 3.55\%. The BBox U-Net and the \gls{dp} seam raise the one-page result to 99.18\%. Without this pretraining, the best initialization reaches 56.23\% with one page on average, compared with 84.53\% for CATMuS, and the selection of the pages has a smaller effect than the initialization.

\textbf{RQ3: How does the diversity of the training collections in script and layout affect the robustness of a few-shot text line segmentation model on unseen manuscript collections?}

The diversity of the source collections matters for Mask R-CNN but not for the center-line model. For Mask R-CNN, this diversity acts mainly through the line styles of the source collections. As described in Section~\ref{sec:disc-transfer-errors}, its transfer decreases as the line sizes of the source and the target collection differ more, and it merges neighboring lines on dense pages. With a fixed budget of six training pages, Mask R-CNN drops from 83.56\% with six source collections to 77.55\% with three on average over three seeds, while the number of source collections changes the result of the center-line model by less than one percentage point. The pretraining alone is robust for the center-line model, which reaches 93.07\% without fine-tuning compared with 49.09\% for Mask R-CNN. This agrees with Islam et al.~\cite{Textline_maskrcnn}, who attribute the transfer errors of Mask R-CNN to a lower interline spacing and not to the script. After training on the six other collections (\gls{loco}-3), both models transfer to unseen collections, with 93.42\% for the center-line model and 85.44\% for Mask R-CNN. Compared with one source collection, Mask R-CNN improves from 70.55\% to 85.44\%, but this comparison also increases the number of training pages.

\section{Future Work}
\label{sec:future-work}

Three directions follow from these findings. First, pretraining on collections with more varied layouts, including densely written pages, may reduce the line merges of the detection-based models. Second, this thesis evaluates the segmentation with line detection metrics only. How the different \gls{tls} pipelines affect the accuracy of a subsequent \gls{htr} model, for example when line masks miss ascenders, descenders, or diacritics, remains to be studied. Third, recent models recognize the text of a complete page without a prior line segmentation~\cite{coquenet2023dan}. Page-level recognition may be more efficient than segmenting and recognizing each line separately, and a comparison of both approaches with few annotated pages remains open.
